\emph{Setting.} By Fact \claimref{scrambling.imported-kostant-madore-partition-data}(i), $\mathfrak{sl}_2$ structures on $\Mat_n$ (Proposition \claimref{scrambling.boundary-mode-space-matrix}) are classified up to conjugation by the partition $\lambda$ of $n$ given by the restriction of the defining representation, $\C^n \cong \bigoplus_d V_d \otimes \C^{m_d}$ with $V_d$ the $d$-dimensional irreducible representation. In a basis adapted to the blocks the standard triple is $e_\lambda$ (one Jordan block of size $d$ per part), $h_\lambda = \bigoplus \operatorname{diag}(d-1, d-3, \dots, 1-d)$ and $f_\lambda$ with entries $(i+1)(d-1-i)$ below the diagonal of each block; then $[h,e] = 2e$, $[h,f] = -2f$, $[e,f] = h$, since the diagonal of $[e,f]$ in a block is $(i+1)(d-1-i) - i(d-i) = d-1-2i$. The check builds these triples exactly for every partition of every $n \leq 6$ and verifies the three relations.

\emph{(i).} The conjugation orbit of $e_\lambda$ is $GL_n/C(e_\lambda)$ with $C(e_\lambda)$ its centraliser, of dimension $n^2 - \dim C(e_\lambda)$. A matrix commuting with $e_\lambda$ is a module map of $\C^n$ over the polynomial ring in $e_\lambda$, and the maps from a Jordan block of size $a$ to one of size $b$ form a space of dimension $\min(a,b)$; so $\dim C(e_\lambda) = \sum_{a,b}\min(\lambda_a, \lambda_b) = \sum_i (\lambda^{T}_i)^2$, the last step counting, for each level $i$, the pairs of parts both exceeding $i$. The check \texttt{check\_principal\_is\_the\_unique\_full\_dimensional\_orbit} computes $\dim C(e_\lambda)$ as the nullity of $X \mapsto [X, e_\lambda]$ by exact rank and confirms the formula for every $\lambda$ with $n \leq 6$. The orbit is a point exactly when $C(e_\lambda) = \mathfrak{gl}_n$, that is when $e_\lambda$ is central and nilpotent, that is $e_\lambda = 0$, $\lambda = (1^n)$; then $h = [e,f] = 0$ and the triple is zero. Conversely an $\mathfrak{sl}_2$ structure $\phi$ fixed pointwise by every conjugation, $\operatorname{Ad}_g \circ \phi = \phi$, has central image, and $\mathfrak{sl}_2 = [\mathfrak{sl}_2, \mathfrak{sl}_2]$ maps to zero in the abelian centre; so the zero map is the only structure that the symmetry of $\Mat_n$ fixes. The check confirms that $(1^n)$ is the only partition with orbit dimension $0$.

\emph{(ii).} By Schur's lemma the commutant of the image of the triple is $\bigoplus_d \mathfrak{gl}(m_d)$, acting on the multiplicity spaces, of dimension $\sum_d m_d^2$; the check computes it as the nullity of $X \mapsto ([X,e], [X,f])$ and confirms the formula for every $\lambda$ with $n \leq 6$. The group $GL(m_d)$ acts transitively on the decompositions of $V_d \otimes \C^{m_d}$ into $m_d$ copies of $V_d$ and contains the permutations of the copies in any one decomposition, so no block of size $d$ is distinguished from another of the same size; an equivariant automorphism preserves dimensions, so blocks of different sizes are always told apart. Hence no block is distinguished exactly when all parts are equal, $\lambda = (k^{n/k})$ with $k \mid n$. The check confirms, for $n \leq 6$, that the rectangular partitions are exactly those indexed by the divisors of $n$, that there are two of them, $(n)$ and $(1^n)$, exactly when $n$ is prime, and that $(2,2)$ at $n = 4$ has a four-dimensional commutant.

\emph{(iii).} $\sum_i (\lambda^{T}_i)^2 \geq \sum_i \lambda^{T}_i = n$, with equality exactly when every $\lambda^{T}_i = 1$, that is $\lambda^{T} = (1^n)$, $\lambda = (n)$; so the principal orbit has dimension $n^2 - n$ and every other orbit is strictly smaller, as the check confirms for $n \leq 6$. Every nilpotent matrix, real or complex, is conjugate to some $e_\lambda$ (its Jordan form; all dimensions here are over the ground field), so the nilpotent cone $\mathcal{N}$ is the finite union of the orbits $\mathcal{O}_\lambda$. The non-principal orbits form the closed set $\{x \in \mathcal{N} : x^{n-1} = 0\}$, so $\mathcal{O}_{(n)}$ is open in $\mathcal{N}$; it is dense, since adding $\varepsilon$ to the missing superdiagonal entries of $e_\lambda$ gives a strictly upper triangular matrix with nonzero superdiagonal, which is regular nilpotent, and tends to $e_\lambda$. The other orbits are submanifolds of dimension less than $n^2 - n$ and so have zero $(n^2 - n)$-dimensional measure; conjugation acts by diffeomorphisms, so the class of smooth positive densities on $\mathcal{O}_{(n)}$ is conjugation invariant, and every measure absolutely continuous with respect to it gives the type $(n)$ with probability one. Through the classification, a nilpotent element of type $(n)$ is the $e$ of the principal embedding.

\emph{Conclusion.} The three readings of indifference give three different answers: the zero map, the rectangular types, the principal embedding. Theorem \claimref{scrambling.principality-indifference-sphere} needs the third, which is a symmetry statement (the invariant measure class) together with a measure principle (a choice made without preference lands where that class puts full measure); the first two do not give it, the second failing for every composite $n$.
